% ============================================================
% ECON 320 — Introduction to Mathematical Economics
% Quiz 1 — Student Version
% Fall 2026  |  Muhebullah Karimzada  |  UNBC
% ============================================================
\documentclass[11pt]{article}
\usepackage[margin=1in,top=1.2in,bottom=1.2in]{geometry}
\usepackage[utf8]{inputenc}
\usepackage[T1]{fontenc}
\usepackage{lmodern}
\usepackage{amsmath,amssymb,bm}
\usepackage{booktabs,array,multirow,adjustbox}
\usepackage{enumitem}
\usepackage{xcolor}
\usepackage{fancyhdr}
\addtolength{\topmargin}{-10pt}
\usepackage{tcolorbox}
\usepackage{titlesec}
\usepackage{hyperref}
\tcbuselibrary{skins}
\hypersetup{colorlinks=false,hidelinks}

%% ── Colours ────────────────────────────────────────────────
\definecolor{UNBCDark}{RGB}{0,0,0}
\definecolor{UNBCGold}{RGB}{0,0,0}
\definecolor{SoftGray}{RGB}{244,246,248}

%% ── Header / Footer ────────────────────────────────────────
\pagestyle{fancy}
\fancyhf{}
\renewcommand{\headrulewidth}{0pt}
\lhead{\small\textbf{Name:\enspace\underline{\hspace{6cm}}}}
\rhead{\small Student ID:\enspace\underline{\hspace{4.5cm}}}
\cfoot{\small\thepage}

%% ── Section style ──────────────────────────────────────────
\titleformat{\section}[block]
  {\large\bfseries\color{UNBCDark}}{}
  {0em}{}
  [{\color{UNBCGold}\titlerule[0.8pt]}]
\titlespacing*{\section}{0pt}{1.2em}{0.6em}

%% ── Custom boxes ───────────────────────────────────────────
\newtcolorbox{infobox}{
  colback=SoftGray, colframe=UNBCDark,
  left=5mm, right=5mm, top=4mm, bottom=4mm,
  boxrule=0.7pt, arc=3mm}

\newtcolorbox{databox}{
  colback=SoftGray, colframe=UNBCDark!40,
  left=4mm, right=4mm, top=3mm, bottom=3mm,
  boxrule=0.4pt, arc=2mm}

% ============================================================
\begin{document}
% ============================================================

%% ── Title block ─────────────────────────────────────────────
\begin{center}
{\Large\bfseries ECON 320 \;--\; Introduction to Mathematical Economics}\\[0.5em]
{\large\bfseries Quiz 1\quad$|$\quad September 24, 2026}\\[0.2em]
{\small University of Northern British Columbia\quad
 Instructor: M.\,Karimzada}
\end{center}

\vspace{0.5em}



%% ════════════════════════════════════════════════════════════
\section{Part A:\enspace Multiple Choice\quad\normalfont\small
         (5 questions $\times$ 1 mark = 5 marks)}
%% ════════════════════════════════════════════════════════════

\medskip

\begin{enumerate}[label=\textbf{Q\arabic*.},leftmargin=2.5em,itemsep=1.6em]

\item Let $U = \{1,2,3,4,5,6,7,8\}$, $A = \{1,2,5,6\}$, and $B = \{2,3,6,7\}$.
Which of the following correctly gives $(A \cap B)^{c}$?

\medskip
\begin{enumerate}[label=\textbf{(\Alph*)},leftmargin=2em,itemsep=0.4em]
  \item $\{2,\,6\}$
  \item $\{1,\,3,\,4,\,5,\,7,\,8\}$
  \item $\{1,\,2,\,3,\,5,\,6,\,7\}$
  \item $\{4,\,8\}$
\end{enumerate}

\item For a set $A$ with $|A| = 5$, the number of elements in the power set
$\mathcal{P}(A)$ is:

\medskip
\begin{enumerate}[label=\textbf{(\Alph*)},leftmargin=2em,itemsep=0.4em]
  \item $10$
  \item $25$
  \item $32$
  \item $64$
\end{enumerate}

\item Which of the following functions $f : \mathbb{R} \to \mathbb{R}$ is
\textbf{injective but not surjective}?

\medskip
\begin{enumerate}[label=\textbf{(\Alph*)},leftmargin=2em,itemsep=0.4em]
  \item $f(x) = x^{2}$
  \item $f(x) = e^{x}$
  \item $f(x) = x^{3}$
  \item $f(x) = \sin x$
\end{enumerate}

\item The inverse of $\displaystyle f(x) = \frac{2x+3}{x-1}$, \; $x \neq 1$, is:

\medskip
\begin{enumerate}[label=\textbf{(\Alph*)},leftmargin=2em,itemsep=0.4em]
  \item $\displaystyle f^{-1}(x) = \frac{x+3}{x-2}$
  \item $\displaystyle f^{-1}(x) = \frac{x-3}{x+2}$
  \item $\displaystyle f^{-1}(x) = \frac{x-1}{2x+3}$
  \item $\displaystyle f^{-1}(x) = \frac{2x-1}{x+3}$
\end{enumerate}

\item Let $f(x) = x^{2} - 1$ and $g(x) = \sqrt{x+2}$. What is $(g \circ f)(3)$?

\medskip
\begin{enumerate}[label=\textbf{(\Alph*)},leftmargin=2em,itemsep=0.4em]
  \item $2\sqrt{2}$
  \item $\sqrt{10}$
  \item $4$
  \item $3$
\end{enumerate}

\end{enumerate}

\bigskip

%% ════════════════════════════════════════════════════════════
\section{Part B:\enspace Numerical Problems\quad\normalfont\small
         (15 marks)}
%% ════════════════════════════════════════════════════════════

\bigskip

%% ── Q6 ───────────────────────────────────────────────────────
\noindent\textbf{Q6.}\quad(5 marks)\enspace
Let $U = \{1,2,3,4,5,6,7,8,9,10\}$, $A = \{1,3,5,7,9\}$,
$B = \{1,2,3,4,5\}$, and $C = \{3,5,6,7,8\}$.

\medskip
\begin{enumerate}[label=\textbf{(\alph*)},leftmargin=2em,itemsep=1.5em]
  \item Find $A \cap B$ and $A \setminus B$.\hfill\textbf{[2 marks]}

        

  \item Find $(A \cup B)^{c}$.\hfill\textbf{[1 mark]}

         

  \item Verify De Morgan's Law: $(A \cup B)^{c} = A^{c} \cap B^{c}$.\hfill\textbf{[1 mark]}

      

  \item Determine $\bigl|\,\mathcal{P}(A \cap B \cap C)\,\bigr|$.\hfill\textbf{[1 mark]}

         
\end{enumerate}

\bigskip

%% ── Q7 ───────────────────────────────────────────────────────
\noindent\textbf{Q7.}\quad(5 marks)\enspace
Let $\displaystyle f(x) = \frac{2x+3}{x-1}$, \; $x \neq 1$.

\medskip
\begin{enumerate}[label=\textbf{(\alph*)},leftmargin=2em,itemsep=1.5em]
  \item Find the natural domain and range of $f$.\hfill\textbf{[2 marks]}

     

  \item Find $f^{-1}(x)$ and state its domain.\hfill\textbf{[2 marks]}

       

  \item Verify that $f\!\left(f^{-1}(x)\right) = x$.\hfill\textbf{[1 mark]}

         
\end{enumerate}

\newpage

%% ── Q8 ───────────────────────────────────────────────────────
\noindent\textbf{Q8.}\quad(5 marks)\enspace
A firm's profit function is $\pi(x) = -2x^{2} + 12x - 10$,
where $x \geq 0$ is the quantity produced.

\medskip
\begin{enumerate}[label=\textbf{(\alph*)},leftmargin=2em,itemsep=1.5em]
  \item Find all quantities at which $\pi(x) = 0$ (the break-even outputs).\hfill\textbf{[1 mark]}

     

  \item Find the profit-maximising quantity $x^{*}$ and the maximum profit $\pi^{*}$.\hfill\textbf{[2 marks]}

        

  \item Is $\pi$ concave or convex? Justify using the second derivative.\hfill\textbf{[1 mark]}

      

  \item Find the range of $\pi$ on $[0,\infty)$.\hfill\textbf{[1 mark]}

       
\end{enumerate}

\vspace{1em}
{\hrule height 0.6pt}
\begin{center}
\small\itshape
--- End of Quiz ---\qquad Total: 20 marks\qquad Time: 20 minutes\qquad Good luck!
\end{center}

\end{document}
